% ============================================================
% 第6章 変数定義と採用根拠：四要素モデルの実証的写像
% ============================================================

\chapter{変数定義と採用根拠：四要素モデルの実証的写像}
\label{chap:variables}

第4章では，Merton モデルおよび Bharath \& Shumway（2008）（以下 B\&S）の
multiplicative → log → linear の構造を基礎として，
新電力企業の倒産構造を反映した潜在倒産距離 $Z_{\text{lin}}$ を理論的に構築した。
本章では，この理論モデルを実証分析に接続するために，
$Z_{\text{lin}}$ の構成要素である収益性 $\mu$，
外部ショック感応度 $\sigma$，支援能力 $S$，倒産閾値 $D$ を，
観測可能な変数として定義し，その採用根拠を整理する。

新電力企業の倒産は，外部ショックが短期間で資金繰りに波及し，
複数の要因が掛け算的に悪化することで急激に顕在化するという特徴を持つ。
この構造は Merton 型モデルの multiplicative な性質と整合的であり，
また B\&S が示した「理論構造を観測可能な proxy によって離散時間に写像する」
という方法論は，新電力企業の倒産リスクを実証的に扱ううえで有効である。

本章では，まず四要素モデルの各構成要素が
Merton/B\&S の理論構造のどの部分に対応するかを明確化し（5.1 節），
続いて収益性 $\mu$（5.2 節），外部ショック感応度 $\sigma$（5.3 節），
支援能力 $S$（5.4 節），倒産閾値 $D$（5.5 節）について，
変数定義と採用根拠を順に整理する。
これにより，第6章で構築する Ordered Logit モデルにおいて，
$Z_{\text{lin}}$ を実証的に推定可能な形で扱うための基盤を整える。

%---------------------------------------------
\section{本章の目的と基本方針}

第4章で構築した潜在倒産距離 $Z_{\text{lin}}$ は，



\[
Z_{\text{lin}} = a\mu - b\sigma + cS - dD
\]



という線形構造を持つ。本章の目的は，この四要素
（収益性 $\mu$，外部ショック感応度 $\sigma$，支援能力 $S$，倒産閾値 $D$）を
実証分析で測定可能な変数として定義し，その採用根拠を明確化することである。

理論モデルを実証に写像する際には，B\&S が行ったように，
観測不可能な理論パラメータを観測可能な proxy によって置き換える必要がある。
本研究では，新電力企業の制度的特徴とデータ制約を踏まえ，
四要素を以下の方針で定義する。

\begin{itemize}
    \item 収益性 $\mu$：総資産の対数差分（log-diff）\footnote{
対数差分（log-diff）とは，$\ln(A_t)-\ln(A_{t-1})$ の形で定義される成長率指標であり，
近似的に $\frac{A_t-A_{t-1}}{A_{t-1}}$（通常の成長率）と一致する。
対数を取ることで，(i) 外れ値の影響を抑え，(ii) 比率変化を加法的に扱うことができ，
(iii) 時系列の安定性が高まるため，経済学・金融研究で標準的に用いられる。
特に財務データのように分布が歪んだ変数では，log-diff の方が統計的性質が良好である。
}
    \item 外部ショック感応度 $\sigma$：スポット価格変化率の標準偏差\footnote{
たとえば，ある 5 日間のスポット価格が
50 円，55 円，60 円，54 円，57 円であったとする。
このときの日次変化率は
$0.10,\ 0.091,\ -0.10,\ 0.055$ となる。
これらの平均は約 $0.036$（3.6\%）であり，
各日の変化率がこの平均からどれだけ離れているかを二乗して平均し，
平方根を取ると標準偏差は約 $0.083$（8.3\%）となる。
この値が「市場ショックの大きさ」を表し，
変化率の揺れが大きいほど $\sigma$ が大きくなり，
新電力企業が受ける収益の不確実性が高まることを意味する。
}
    \item 支援能力 $S$：支援の有無を示す dummy（本研究では $S=0$）
    \item 倒産閾値 $D$：負債比率（Liabilities / Assets）
\end{itemize}

これらの変数は，第6章の Ordered Logit 推定において，
倒産ステージ（存続・財務悪化・倒産）を説明する主要因子として用いる。

%---------------------------------------------
\section{収益性 $\mu$ の定義と採用根拠}

\subsection{理論的対応：企業価値の期待成長率としての $\mu$}

Merton モデルにおける収益性 $\mu$ は企業価値 $V_t$ の期待成長率であり，
倒産距離 $Z_{\text{lin}}$ の正方向の要因として機能する。
企業価値が高い成長率を維持していれば，企業価値が負債額 $F$ を下回る確率は低くなり，
倒産距離は延伸する。

企業金融論では，企業価値の成長率は以下の要因によって規定される。

\begin{itemize}
    \item 投資活動（資本ストックの拡大）
    \item 営業活動（キャッシュフローの創出）
    \item 財務活動（資本構成の改善）
\end{itemize}

しかし新電力企業は非上場であり，企業価値 $V_t$ を市場価格から観測することはできない。
また，損益計算書が欠損する企業が多く，営業利益率や ROA を安定的に計算することも困難である。

このため，企業価値の期待成長率を観測可能な財務指標で代理する必要がある。

\subsection{総資産の対数差分（log-diff）による実証的定義}

本研究では，企業価値の期待成長率を総資産の対数差分（log-diff）によって代理する。



\[
\mu_{i,t}
= \ln(\mathrm{Assets}_{i,t}) - \ln(\mathrm{Assets}_{i,t-1})
\]



総資産は企業価値の蓄積を反映するストック指標であり，
投資活動・営業活動・財務活動の結果がすべて総資産に集約される。
したがって，総資産の成長率は企業価値の期待成長率の proxy として理論的に妥当である。

log-diff を用いる理由は以下の通りである。

\begin{itemize}
    \item 成長率を加法的に扱えるため，線形モデル（B\&S の構造）と整合する
    \item 外れ値の影響を抑え，非正規性に強い
    \item 資産規模の違いによるスケール問題を回避できる
\end{itemize}

\subsection{企業金融論の視点：総資産成長率が倒産距離を延伸する理由}

企業金融論では，総資産の成長率が倒産リスクを低減するメカニズムは以下の通りである。

\begin{itemize}
    \item \textbf{資本ストックの拡大}：
          投資活動による資産増加は企業価値の増加を意味し，
          倒産閾値 $D$ からの距離を広げる。
    \item \textbf{営業活動の改善}：
          営業キャッシュフローが資産として蓄積されることで，
          資金繰りの安定性が向上する。
    \item \textbf{財務活動の強化}：
          増資や内部留保の蓄積は自己資本を増加させ，
          レバレッジを低下させることで倒産距離を延伸する。
    \item \textbf{ショック吸収能力の向上}：
          資産規模が大きい企業ほど，外部ショック（価格急騰）を吸収しやすい。
\end{itemize}

これらの構造は，新電力企業の倒産メカニズムと完全に整合する。

\subsection{新電力企業の産業特性：総資産が企業価値の主要構成要素となる理由}

新電力企業は以下の産業特性を持つ。

\begin{itemize}
    \item 発電設備を保有しない「アセットライト型」企業が多い
    \item 営業利益率が価格ショックに強く依存し，安定的に観測できない
    \item 資産の大部分が流動資産（売掛金・預金）であり，資金繰りに直結する
    \item 投資活動よりも営業活動・財務活動の影響が総資産に強く反映される
\end{itemize}

このため，総資産の成長率は企業価値の期待成長率を最も安定的に捉える指標となる。

\subsection{英国企業における補助指標：総資産成長率の質的側面}

英国企業では，財務諸表の開示制度が整備されており，
総資産の成長率の「質的側面」を補完する指標が観測可能である。

\begin{itemize}
    \item \textbf{固定資産の構成}：発電設備・ITインフラへの投資状況
    \item \textbf{流動資産の構成}：売掛金・預金の増加による資金繰り改善
    \item \textbf{資本構成の変化}：増資・内部留保の蓄積による自己資本の増加
    \item \textbf{関連会社取引}：親会社支援による資産増加（intercompany loan）
\end{itemize}

これらの補助指標は，総資産成長率の背後にある企業価値の変化をより精緻に把握するために有用である。

\subsection{採用根拠：理論・制度・データの三側面からの整合性}

\begin{itemize}
    \item Merton モデルの企業価値成長率 $\mu$ に直接対応する
    \item 総資産は企業価値の蓄積を反映するストック指標であり，非上場企業でも安定的に観測可能
    \item 新電力企業の産業特性（アセットライト型・ショック依存型）と整合する
    \item 英国企業では総資産の質的側面を補完する情報が豊富である
    \item Ordered Logit モデル（第7章・第8章）において，
          倒産ステージの説明変数として理論的・実証的に妥当である
\end{itemize}

以上より，総資産の対数差分は，
理論構造・企業金融論・産業特性・データ制約のいずれにも整合的な収益性指標である。


%---------------------------------------------
\section{外部ショック感応度 $\sigma$ の定義と採用根拠}

\subsection{理論的対応：Merton モデルにおける不確実性の役割}

Merton モデルにおける $\sigma_V$ は企業価値のボラティリティであり，
倒産距離の分母として機能する。すなわち，$\sigma_V$ が大きいほど
企業価値の変動幅が広がり，企業価値が負債額 $F$ を下回る確率が高まる。

企業金融論では，ボラティリティは倒産リスクを規定する中心的要因であり，
特にレバレッジの高い企業では，価値変動が自己資本に与える影響が増幅される。
この構造は，新電力企業の倒産メカニズムと整合的である。

新電力企業は，外部ショック（スポット価格急騰・インバランス単価上昇）が
短期間で資金繰りに波及する「ショック依存型産業」であり，
企業価値のボラティリティは市場価格の変動に強く依存する。
したがって，$\sigma_V$ を市場価格の変動から proxy 化することは理論的に妥当である。

\subsection{新電力企業の倒産構造：ショック依存型ボラティリティ}

新電力企業の倒産構造には，以下の特徴がある。

\begin{itemize}
    \item \textbf{スポット市場依存度が高い}：調達費用の大部分がスポット価格に連動する
    \item \textbf{ショックが即時に損益へ波及}：価格急騰がそのまま資金繰り悪化につながる
    \item \textbf{ショックが負債構造を内生的に変化させる}：
          インバランス支払や調達費用が急増し，負債額が短期間で膨張する
    \item \textbf{ショックが支援能力と相互作用する}：
          親会社支援がある企業はショックを吸収できるが，支援能力が低い企業は吸収できない
\end{itemize}

これらの特徴は，外部ショック感応度 $\sigma$ が倒産距離の分母として
極めて重要な役割を果たすことを示している。

\subsection{実証上の定義：価格変化率の標準偏差}

本研究では，市場ショックの強さをスポット価格の変化率の標準偏差として定義する。



\[
r_t = \frac{P_t - P_{t-1}}{P_{t-1}}, \qquad
\sigma = \sqrt{\frac{1}{N-1}\sum_{t=1}^{N}(r_t - \bar{r})^2}
\]



変化率を用いることで価格水準の違いによるスケール問題を回避でき，
市場ショックの強度を適切に捉えることができる。

\subsection{企業金融論の視点：ボラティリティが倒産リスクを規定する理由}

企業金融論では，ボラティリティが倒産リスクを規定するメカニズムは以下の通りである。

\begin{itemize}
    \item \textbf{価値変動の増幅}：
          レバレッジの高い企業では，価値変動が自己資本に与える影響が増幅される。
    \item \textbf{資金繰り制約の悪化}：
          ショックが短期負債の増加を引き起こし，資金繰りが急速に悪化する。
    \item \textbf{倒産閾値との相互作用}：
          ショックが負債額を押し上げ，倒産閾値 $D$ を内生的に引き上げる。
    \item \textbf{支援能力との相互作用}：
          支援能力 $S$ が高い企業はショックを吸収できるが，
          $S$ が低い企業はショックがそのまま倒産リスクに直結する。
\end{itemize}

これらの構造は，新電力企業の倒産メカニズムと完全に整合する。

\subsection{英国企業におけるショック感応度の補助指標}

英国企業では，財務諸表・登記情報から以下の補助指標を取得できる。

\begin{itemize}
    \item \textbf{調達構造の開示}：固定価格契約の比率，スポット依存度
    \item \textbf{担保付負債の有無}：ショック時の資金調達能力を反映
    \item \textbf{関連会社貸付（intercompany loan）}：
          ショック吸収能力（内部資本市場の厚み）を反映
    \item \textbf{親会社の財務体力}：ショック時の支援能力と相互作用
\end{itemize}

これらはショック感応度の「質的側面」を補完する指標であり，
英国企業を対象とする場合には，価格変化率の標準偏差と併用することで
ショック感応度の構造をより精緻に捉えることができる。

\subsection{採用根拠：理論・制度・データの三側面からの整合性}

\begin{itemize}
    \item Merton モデルの $\sigma_V$ に対応し，倒産距離の分母として機能する
    \item 新電力企業の倒産構造（ショック依存型産業）と完全に整合する
    \item 価格変化率の標準偏差は市場ショックの強度を最も適切に捉える
    \item 英国企業では調達構造・担保設定・内部資本市場などの補助指標が利用可能である
    \item 支援能力 $S$ や倒産閾値 $D$ と相互作用し，Ordered Logit モデルの説明力を高める
\end{itemize}

以上より，外部ショック感応度 $\sigma$ は，
新電力企業の倒産構造を反映する主要因子として理論的・実証的に妥当である。

%---------------------------------------------

\section{支援能力 $S$ の定義と採用根拠}

\subsection{理論的対応：資本構成と倒産距離の関係}

支援能力 $S$ は Merton/B\&S には存在しない，本研究独自の拡張要素である。
企業金融論において，資本構成は倒産リスクを規定する中心的概念であり，
親会社・金融機関・自治体など外部主体による支援は，
実質的に倒産閾値を引き下げ，倒産距離を延伸する効果を持つ。

Merton モデルでは，企業価値 $V_t$ が負債額 $F$ を下回ると倒産が発生する。
しかし新電力企業の倒産構造では，外部ショックにより資金繰りが急速に悪化した際，
親会社や金融機関が資金注入・信用補完・債務肩代わりを行うことで，
企業価値が一時的に $F$ を下回っても倒産が発生しないケースが存在する。

この構造は，倒産閾値 $D$ が外部主体によって「内生的に引き下げられる」ことを意味し，
支援能力 $S$ は倒産距離 $Z_{\text{lin}}$ の正方向の要因として機能する。



\[
Z_{\text{lin}} = a\mu - b\sigma + cS - dD
\]



企業金融論の観点からは，$S$ は
「資本構成の柔軟性」「資金調達能力」「グループ内部資本市場の厚み」
を反映する変数として解釈できる。

\subsection{日本企業における観測制約}

日本の新電力企業では，財務諸表・登記情報から
親会社支援や信用補完の客観的証拠を読み取ることが困難である。
増資・関連会社貸付・担保設定・役員派遣などの支援行動が
財務諸表に明示されない場合が多く，支援能力を連続変数として測定することはできない。

そのため，日本企業を対象とした分析では，
支援能力を二値の dummy として定義し，分析対象企業では
支援の客観的証拠が確認できなかったため $S_i = 0$ とした。

\subsection{英国企業における支援能力の観測可能性：資本構成情報の豊富さ}

英国企業では Companies House における登記制度が充実しており，
親会社支援の「痕跡」を複数の公式記録から観測できる。
これは企業金融論の観点から極めて重要であり，
資本構成の柔軟性・資金調達能力・グループ内部資本市場の厚みを
実証的に測定することが可能となる。

具体的には以下の情報が支援能力の proxy として利用可能である。

\begin{itemize}
    \item \textbf{増資（share allotment）}：
          親会社が危機時に資本注入を行ったか。  
          これは資本構成の強化（Equity buffer の増加）を直接反映する。
    \item \textbf{関連会社貸付（intercompany loan）}：
          親会社が lender として資金支援を行ったか。  
          グループ内部資本市場の厚みを示す。
    \item \textbf{担保設定（charge）}：
          親会社や大手銀行が信用補完を提供しているか。  
          負債の実効コストを引き下げ，倒産閾値を内生的に低下させる。
    \item \textbf{支配株主（Persons with Significant Control）}：
          親会社の財務体力・資本構成を反映し，支援主体の能力を示す。
    \item \textbf{役員派遣（director appointment）}：
          親会社がガバナンス支援を行っているか。  
          経営資源の投入は支援能力の重要なシグナルである。
    \item \textbf{企業グループ構造（group structure）}：
          グループ内での資金移転可能性・内部資本市場の存在を示す。
\end{itemize}

これらは日本企業では観測できないが，
英国企業では公式記録として入手可能であり，
支援能力 $S$ を連続的な潜在変数として構築することが可能となる。

\subsection{実証上の定義：資本構成の潜在因子としての $S$}

本研究では，日本企業と英国企業の双方を対象とするため，
支援能力 $S$ を以下のように定義する。



\[
S_i = f(\text{増資}, \text{関連会社貸付}, \text{担保設定}, 
       \text{支配株主}, \text{役員派遣}, \text{グループ構造})
\]



企業金融論の観点からは，これらの proxy は
「資本構成の柔軟性」「資金調達能力」「内部資本市場の厚み」
を反映する変数群であり，支援能力の潜在因子として統合することが合理的である。

具体的には，これらの proxy 変数を因子分析（Factor Analysis）または
主成分分析（PCA）によって 1 次元の潜在変数として抽出し，
支援能力 $S$ を連続変数として扱う。

日本企業については支援の痕跡が観測できないため $S_i = 0$ とし，
英国企業については上記の proxy から推定した $S_i$ を用いる。

\subsection{採用根拠：企業金融論との整合性}

\begin{itemize}
    \item 支援能力は資本構成の柔軟性を反映し，
          倒産閾値を内生的に引き下げる理論的役割を持つ
    \item 英国企業では支援能力の痕跡が公式記録として観測可能であり，
          日本企業よりも S の実証的測定が可能である
    \item 増資・貸付・担保設定・役員派遣などの支援行動は，
          資本構成の改善・資金調達能力の強化・信用補完の提供という
          企業金融論の主要メカニズムと整合する
    \item PCA による潜在変数化により，複数の proxy を統合し，
          支援能力の強度を連続的に測定できる
\end{itemize}

以上より，支援能力 $S$ は日本企業では dummy として扱い，
英国企業では資本構成情報に基づく潜在変数として推定することで，
理論構造とデータ制約の双方に整合的な実証モデルを構築できる。


%---------------------------------------------
\section{倒産閾値 $D$ の定義と採用根拠}

\subsection{理論的対応：Merton モデルにおける負債額 $F$ の役割}

Merton モデルでは，倒産は企業価値 $V_t$ が負債額 $F$ を下回った時点で発生する。
すなわち $F$ は倒産閾値として機能し，企業金融論の観点では
「資本構成（Leverage）」「負債の性質」「債務の優先順位」
を反映する中心的パラメータである。

負債額が大きいほど倒産閾値は押し上げられ，倒産距離は短縮する。
この構造は，資本構成が倒産リスクを規定するという企業金融論の基本命題と整合する。

しかし新電力企業の倒産は，固定負債の返済不能ではなく，
外部ショック（インバランス単価急騰・調達費用高騰）が
短期間で資金繰りに波及することで発生する。
したがって，Merton モデルの負債額 $F$ をそのまま用いることはできず，
外部ショック時の支払義務を proxy によって実証的に写像する必要がある。

\subsection{新電力企業の倒産構造：負債の「性質」の違い}

新電力企業の倒産構造には，以下の特徴がある。

\begin{itemize}
    \item 負債の大部分が短期流動負債（仕入債務・インバランス支払）である
    \item 外部ショック時に支払義務が急増し，負債額が内生的に変動する
    \item 固定負債よりも「支払義務の急増」が倒産の直接原因となる
    \item 資金繰り悪化が倒産のトリガーであり，財務レバレッジの影響が強い
\end{itemize}

このため，倒産閾値を「固定負債額」として扱うことは不適切であり，
負債の水準と構成を反映する財務指標を用いる必要がある。

\subsection{実証上の定義：負債比率による倒産閾値の proxy}

外部ショック時の支払義務は企業別に観測できないため，
本研究では財務指標である負債比率を倒産閾値の proxy として用いる。



\[
D_{i,t} = \frac{\mathrm{Liabilities}_{i,t}}{\mathrm{Assets}_{i,t}}
\]



負債比率は，資本構成の中心指標である Leverage を直接反映し，
倒産閾値の押し上げ効果を測定するうえで最も安定した proxy である。

\subsection{企業金融論の視点：負債比率が倒産閾値を規定する理由}

企業金融論では，負債比率が倒産リスクを規定するメカニズムは以下の通りである。

\begin{itemize}
    \item \textbf{レバレッジ効果}：
          負債比率が高いほど，企業価値の下落が自己資本に与える影響が大きく，
          倒産閾値に早く到達する。
    \item \textbf{債務の優先順位}：
          流動負債が多い企業ほど，外部ショック時に支払義務が即時に発生し，
          倒産閾値が事実上引き上げられる。
    \item \textbf{資金繰り制約}：
          新電力企業は短期債務依存度が高く，資金繰り悪化が倒産の直接原因となる。
    \item \textbf{負債の性質}：
          インバランス支払やスポット調達費用は「ショック依存型負債」であり，
          価格変動に応じて負債額が内生的に変動する。
\end{itemize}

これらの特徴は，負債比率が倒産閾値の proxy として最も適切であることを示している。

\subsection{英国企業における倒産閾値の観測可能性}

英国企業では，財務諸表の開示制度が整備されており，
負債の構成（流動負債・長期負債），担保設定（charge），
債務の優先順位，関連会社貸付などの情報が詳細に観測できる。

これにより，倒産閾値の構造をより精緻に把握できる。

\begin{itemize}
    \item 流動負債比率（Current Liabilities / Assets）
    \item 担保付負債の有無（secured debt）
    \item 親会社による信用補完（guarantee）
    \item 関連会社貸付（intercompany loan）
\end{itemize}

これらは倒産閾値の「質的側面」を反映する指標であり，
英国企業を対象とする場合には，負債比率と併用することで
倒産閾値の構造をより正確に捉えることができる。

\subsection{採用根拠：理論・制度・データの三側面からの整合性}

\begin{itemize}
    \item Merton モデルの負債額 $F$ に対応し，倒産閾値の押し上げ効果を直接反映する
    \item 新電力企業の倒産構造（ショック依存型負債・資金繰り悪化）と整合する
    \item 財務諸表から安定的に観測可能であり，企業間比較が容易である
    \item 英国企業では負債の構成・担保・信用補完などの補助情報も利用可能である
    \item Ordered Logit モデル（第7章・第8章）において，
          倒産ステージの説明変数として理論的・実証的に妥当である
\end{itemize}

以上より，倒産閾値 $D$ は負債比率を中心指標としつつ，
英国企業では負債構成・担保設定・信用補完などの補助指標を併用することで，
理論構造と企業金融論の視点の双方に整合的な実証モデルを構築できる。

%---------------------------------------------
\section{本章のまとめ}

本章では，理論モデル $Z_{\text{lin}}$ の四要素
（$\mu, \sigma, S, D$）を実証分析に接続するための
観測可能な proxy として定義し，その採用根拠を整理した。

\begin{itemize}
    \item $\mu$：総資産の対数差分（企業価値の期待成長）
    \item $\sigma$：スポット価格変化率の標準偏差（市場ショックの強度）
    \item $S$：支援能力 dummy（外部支援の有無）
    \item $D$：負債比率（倒産閾値の押し上げ効果）
\end{itemize}

これらの変数は，第6章で構築する Ordered Logit モデルにおいて，
倒産ステージ（存続・財務悪化・倒産）を説明する主要因子として用いる。
本章は，理論モデル（第4章）と実証モデル（第6章）をつなぐ橋渡しとして位置づけられる。
